
\documentclass{article}
\usepackage{fontspec}
\setmainfont{Times New Roman}
\setmonofont{Menlo}

\usepackage{amsmath,amssymb,booktabs,longtable,array,calc,graphicx,xcolor,hyperref}
\usepackage[margin=25mm]{geometry}
\hypersetup{colorlinks=true,linkcolor=blue,urlcolor=blue}
\setlength{\parindent}{0pt}
\setlength{\parskip}{6pt}
\providecommand{\tightlist}{\setlength{\itemsep}{0pt}\setlength{\parskip}{0pt}}
\providecommand{\pandocbounded}[1]{\begin{center}\resizebox{0.96\linewidth}{!}{#1}\end{center}}
\newcounter{none}
\makeatletter
\define@key{Gin}{alt}{}
\makeatother

\begin{document}
\section{Experimental Specification for the White-Cat, Black-Cat, and
Gray-Cat Metastable Interface in AB-C-D Stages
v1}\label{experimental-specification-for-the-white-cat-black-cat-and-gray-cat-metastable-interface-in-ab-c-d-stages-v1}

\textbf{Subtitle:} Closed-system numerical experiment on the AB two-body
metastable interface, weak C readout, and strong D observation for
white-cat, black-cat, and gray-cat branching\\
\textbf{Date:} 2026-07-14\\
\textbf{Author:} Noriaki Kihara\\
\textbf{Series:} Wave Information Readout / staged experiment for
macroscopic cat-like metastable states\\
\textbf{Target implementation:} Python numerical experiment\\
\textbf{Version DOI:} 10.5281/zenodo.21353209\\
\textbf{Concept DOI:} 10.5281/zenodo.21353208

\begin{center}\rule{0.5\linewidth}{0.5pt}\end{center}

\subsection{0. Summary}\label{summary}

This experiment does not assume in advance that an equal-allocation
state of the white-cat state \texttt{A} and the black-cat state
\texttt{B} is merely an unselected superposition.

First, the AB two-body system alone is used to distinguish three phases.

\begin{verbatim}
gray eigen phase:
  A/B remains fixed near 0.5/0.5.

gray metastable phase:
  A/B oscillates weakly near 0.5/0.5 but does not naturally converge to one side.

natural selection phase:
  A or B is selected even without C or D.
\end{verbatim}

Then C is added to test whether the metastable state can be read without
selection. Finally D is added to test whether observation selects either
A or B.

The first checkpoint is not the action of C or D.

\begin{verbatim}
The first checkpoint is to find the interface between the gray eigen phase,
gray metastable phase, and natural selection phase in AB alone.
\end{verbatim}

\begin{center}\rule{0.5\linewidth}{0.5pt}\end{center}

\subsection{1. Purpose}\label{purpose}

The experiment checks the following stages in order.

\begin{enumerate}
\def\labelenumi{\arabic{enumi}.}
\tightlist
\item
  Classify the gray eigen phase, gray metastable phase, and natural
  selection phase in AB alone.
\item
  Add C and test whether a gray eigen or metastable state can be read
  without selection.
\item
  Add D and test the conditions under which observation selects A or B.
\item
  Confirm whether a gray eigen phase remains gray even under D.
\item
  In large-amplitude regions where A/B is already separated, confirm
  that D agrees with the C readout.
\end{enumerate}

D is not assumed to always cause collapse.

If a gray cat has already become an eigen state, it does not need to
fall into white or black under D.

If A/B is already clearly separated, the D result is treated as a
readout of an already existing A/B branching, not as a selection created
by D.

\begin{center}\rule{0.5\linewidth}{0.5pt}\end{center}

\subsection{2. State Variables}\label{state-variables}

At minimum, the following variables are recorded.

{\def\LTcaptype{none} % do not increment counter
\begin{longtable}[]{@{}
  >{\raggedright\arraybackslash}p{(\linewidth - 2\tabcolsep) * \real{0.5000}}
  >{\raggedright\arraybackslash}p{(\linewidth - 2\tabcolsep) * \real{0.5000}}@{}}
\toprule\noalign{}
\begin{minipage}[b]{\linewidth}\raggedright
variable
\end{minipage} & \begin{minipage}[b]{\linewidth}\raggedright
meaning
\end{minipage} \\
\midrule\noalign{}
\endhead
\bottomrule\noalign{}
\endlastfoot
\texttt{a} & complex amplitude of white-cat A \\
\texttt{b} & complex amplitude of black-cat B \\
\texttt{p\_A\ =\ \textbar{}a\textbar{}\^{}2} & white-cat allocation \\
\texttt{p\_B\ =\ \textbar{}b\textbar{}\^{}2} & black-cat allocation \\
\texttt{Q\ =\ p\_A\ +\ p\_B} & total A/B amount \\
\texttt{S\ =\ p\_A\ -\ p\_B} & A/B selection order variable \\
\texttt{S\_amp} & oscillation amplitude of \texttt{S} \\
\texttt{S\_drift} & long-time drift of \texttt{S} \\
\texttt{C\_A,\ C\_B} & A/B readout by C \\
\texttt{D\_A,\ D\_B} & A/B readout by D \\
\texttt{L\_A,\ L\_B} & localization of A/B \\
\texttt{N\_eff\_A,\ N\_eff\_B} & effective harmonic order of A/B \\
\texttt{E\_total} & total conserved quantity or conservation registry
quantity \\
\end{longtable}
}

The central checks are:

\begin{verbatim}
whether Q is conserved;
whether S is fixed, weakly oscillating, or naturally growing;
whether C and D change the phase classification of S.
\end{verbatim}

\begin{center}\rule{0.5\linewidth}{0.5pt}\end{center}

\subsection{3. AB Two-Body Metastable Interface
Experiment}\label{ab-two-body-metastable-interface-experiment}

\subsection{3.1 Purpose}\label{purpose-1}

The AB system alone is used to test what kind of state the gray cat is.

\begin{verbatim}
C = off
D = off
\end{verbatim}

No observation-induced selection is included at this stage.

\subsection{3.2 Initial State}\label{initial-state}

The reference initial values are

\[
a_0 = \frac{1}{\sqrt{2}},
\]

\[
b_0 = \frac{e^{i\phi_0}}{\sqrt{2}}.
\]

Then

\[
Q_0 = |a_0|^2 + |b_0|^2 = 1,
\]

\[
S_0 = |a_0|^2 - |b_0|^2 = 0.
\]

When a small initial bias is inserted,

\[
p_{A,0} = \frac12 + s_0,
\]

\[
p_{B,0} = \frac12 - s_0.
\]

\subsection{3.3 AB Exchange Map}\label{ab-exchange-map}

The minimal AB exchange map is

\[
\begin{pmatrix}
a_{k+1}\\
b_{k+1}
\end{pmatrix}
=
U_\epsilon
\begin{pmatrix}
a_k\\
b_k
\end{pmatrix},
\]

\[
U_\epsilon
=
\begin{pmatrix}
\cos\epsilon & i\sin\epsilon\\
i\sin\epsilon & \cos\epsilon
\end{pmatrix}.
\]

This map alone gives a simple unitary exchange oscillation.

If needed, a weak restoring term or weak nonlinear term is added to test
closed-system stability. However, a term that causes one-sided natural
convergence without C or D is not used for the cat-selection experiment.

\subsection{3.4 Sweep Parameters}\label{sweep-parameters}

At minimum, the following parameters are swept.

{\def\LTcaptype{none} % do not increment counter
\begin{longtable}[]{@{}ll@{}}
\toprule\noalign{}
parameter & meaning \\
\midrule\noalign{}
\endhead
\bottomrule\noalign{}
\endlastfoot
\texttt{epsilon} & strength of AB exchange oscillation \\
\texttt{phi\_0} & initial relative phase \\
\texttt{s\_0} & small initial A/B bias \\
\texttt{noise\_amp} & numerical fluctuation or small disturbance \\
\texttt{K} & number of evolution steps \\
\end{longtable}
}

When the gray cat is treated as a harmonic metastable state, the
amplitude of \texttt{S} is kept small.

The target condition is

\begin{verbatim}
S_amp < S_gray_limit.
\end{verbatim}

The initial guide values are:

\begin{verbatim}
S_gray_limit = 0.05
S_amp target = about 0.01 to 0.03
\end{verbatim}

\subsection{3.5 Phase Classification}\label{phase-classification}

The AB two-body results are classified into three phases.

\subsubsection{Gray Eigen Phase}\label{gray-eigen-phase}

\begin{verbatim}
S_mean ≈ 0
S_amp -> 0
S_drift ≈ 0
Q ≈ 1
\end{verbatim}

In this case, the equal A/B allocation is treated as a new gray-cat
eigen state.

\subsubsection{Gray Metastable Phase}\label{gray-metastable-phase}

\begin{verbatim}
S_mean ≈ 0
0 < S_amp < S_gray_limit
S_drift ≈ 0
Q ≈ 1
\end{verbatim}

This phase is the main target for the D selection experiment.

\subsubsection{Natural Selection Phase}\label{natural-selection-phase}

\begin{verbatim}
abs(S) grows without C or D
S -> +1 or S -> -1
\end{verbatim}

In this phase, observation-induced selection cannot be claimed.

\begin{center}\rule{0.5\linewidth}{0.5pt}\end{center}

\subsection{4. C Readout Experiment}\label{c-readout-experiment}

\subsection{4.1 Purpose}\label{purpose-2}

C is added to a gray eigen or gray metastable state found in AB.

\begin{verbatim}
C = on
D = off
\end{verbatim}

C reads the A/B allocation. At this stage, C must not push the system
into one side.

\subsection{4.2 C Readout Quantities}\label{c-readout-quantities}

C reads:

\[
C_A \approx p_A,
\]

\[
C_B \approx p_B.
\]

The error is recorded as

\[
E_C = |C_A-p_A| + |C_B-p_B|.
\]

\subsection{4.3 C Acceptance Conditions}\label{c-acceptance-conditions}

For the gray metastable phase:

\begin{verbatim}
E_C < tol_C
abs(S_afterC - S_beforeC) < epsilon_C
no selection
\end{verbatim}

For the gray eigen phase:

\begin{verbatim}
C_A ≈ 0.5
C_B ≈ 0.5
S remains near 0
\end{verbatim}

For the large-amplitude region, C reads the current A/B bias.

\begin{verbatim}
C_A > C_B if S > 0
C_B > C_A if S < 0
\end{verbatim}

\begin{center}\rule{0.5\linewidth}{0.5pt}\end{center}

\subsection{5. D Observation Experiment}\label{d-observation-experiment}

\subsection{5.1 Purpose}\label{purpose-3}

D is then added.

\begin{verbatim}
C = optional
D = on
\end{verbatim}

D is a strong observation map used to test whether the system falls into
either A or B.

The result is classified by state phase.

\subsection{5.2 D in the Gray Eigen
Phase}\label{d-in-the-gray-eigen-phase}

For the gray eigen phase, the test asks whether the gray state remains
gray under D.

The expected judgment is:

\begin{verbatim}
D_A ≈ 0.5
D_B ≈ 0.5
S remains near 0
\end{verbatim}

In this case, not falling into white or black is the result.

\subsection{5.3 D in the Gray Metastable
Phase}\label{d-in-the-gray-metastable-phase}

For the gray metastable phase, the test asks whether D selects one side.

The judgment is:

\begin{verbatim}
D_A ≈ 1, D_B ≈ 0
\end{verbatim}

or

\begin{verbatim}
D_A ≈ 0, D_B ≈ 1
\end{verbatim}

If the same fall occurs without D, it is not judged as D-induced
selection.

\subsection{5.4 D in the Large-Amplitude Separated
Region}\label{d-in-the-large-amplitude-separated-region}

When A/B is already clearly separated, D only needs to read the current
A/B bias.

The D readout should agree with the C readout:

\begin{verbatim}
sign(D_A - D_B) = sign(C_A - C_B)
\end{verbatim}

If this agreement is obtained, D is judged to have read an already
existing A/B bias, rather than creating the selection.

\begin{center}\rule{0.5\linewidth}{0.5pt}\end{center}

\subsection{6. Four-Stage Main
Experiment}\label{four-stage-main-experiment}

\subsection{Stage 1: AB Metastable Interface
Search}\label{stage-1-ab-metastable-interface-search}

\begin{verbatim}
C = off
D = off
\end{verbatim}

Output:

\begin{verbatim}
gray eigen phase
gray metastable phase
natural selection phase
\end{verbatim}

\subsection{Stage 2: C Readout
Confirmation}\label{stage-2-c-readout-confirmation}

\begin{verbatim}
C = on
D = off
\end{verbatim}

Output:

\begin{verbatim}
C_read_error
C_induced_bias
C_selection_triggered
\end{verbatim}

\subsection{Stage 3: D Observation
Confirmation}\label{stage-3-d-observation-confirmation}

\begin{verbatim}
C = off or recorded
D = on
\end{verbatim}

Output:

\begin{verbatim}
D_result
D_induced_selection
D_vs_C_agreement
\end{verbatim}

\subsection{Stage 4: C-D Comparison}\label{stage-4-c-d-comparison}

The same AB initial condition is compared under C readout and D
observation.

Checks:

\begin{verbatim}
gray eigen phase:
  C and D both read 0.5/0.5.

gray metastable phase:
  C reads but does not select.
  D is tested for one-sided selection.

large-amplitude separated region:
  C and D read the same A/B bias.
  This is not treated as observation-created selection.
\end{verbatim}

\begin{center}\rule{0.5\linewidth}{0.5pt}\end{center}

\subsection{7. Acceptance Conditions}\label{acceptance-conditions}

The minimum acceptance conditions are:

\begin{enumerate}
\def\labelenumi{\arabic{enumi}.}
\tightlist
\item
  AB alone can classify the gray eigen phase, gray metastable phase, and
  natural selection phase.
\item
  In the gray metastable phase, C can read the A/B allocation without
  selecting one side.
\item
  In the gray eigen phase, D does not make the system fall into white or
  black.
\item
  In the gray metastable phase, D-induced selection can be judged by
  comparison with a no-D control.
\item
  In the large-amplitude separated region, D agrees with the C readout.
\item
  \texttt{Q} or the total conservation registry remains conserved within
  tolerance.
\end{enumerate}

\begin{center}\rule{0.5\linewidth}{0.5pt}\end{center}

\subsection{8. Failure Conditions}\label{failure-conditions}

The experiment does not proceed to later C/D stages if:

\begin{enumerate}
\def\labelenumi{\arabic{enumi}.}
\tightlist
\item
  All AB-only conditions naturally fall into selection.
\item
  \texttt{Q} is not conserved.
\item
  C always causes selection.
\item
  C cannot read the A/B allocation.
\item
  D/no-D differences cannot be classified.
\end{enumerate}

\begin{center}\rule{0.5\linewidth}{0.5pt}\end{center}

\subsection{9. Result Fields}\label{result-fields}

\subsection{9.1 Stage 1: AB Metastable Interface
Search}\label{stage-1-ab-metastable-interface-search-1}

Output:

\begin{verbatim}
gray_cat_ab_metastable_interface_preliminary_result_v1/
\end{verbatim}

Result:

\begin{verbatim}
total_cases = 5600
gray_eigen = 1248
gray_metastable = 733
large_oscillation = 470
natural_selection = 1070
unstable_or_drifting = 2079
\end{verbatim}

\subsection{9.2 Stage 2: C Readout
Window}\label{stage-2-c-readout-window}

Output:

\begin{verbatim}
gray_cat_c_readout_window_preliminary_result_v1/
\end{verbatim}

Result:

\begin{verbatim}
total_cases = 1080
C_window_count = 247
C_informative_window_count = 144
C_nonzero_backaction_window_count = 114
\end{verbatim}

\subsection{9.3 Stage 3: D Observation
Response}\label{stage-3-d-observation-response}

Output:

\begin{verbatim}
gray_cat_d_observation_response_preliminary_result_v1/
\end{verbatim}

Result:

\begin{verbatim}
total_cases = 7056
D_induced_selection_count = 2016
gray_kept_eigen = 1062
white_selected = 1244
black_selected = 772
\end{verbatim}

\subsection{9.4 Stage 4: D Selection
Boundary}\label{stage-4-d-selection-boundary}

Output:

\begin{verbatim}
gray_cat_d_selection_boundary_preliminary_result_v1/
\end{verbatim}

Result:

\begin{verbatim}
boundary_points = 438
selection_possible_boundary_points = 438
no_selection_boundary_points = 0
min_D_gain_overall = 0.0225
max_min_D_gain_overall = 0.065
\end{verbatim}
\end{document}
